\documentclass[../candidatura.tex]{subfiles}
\graphicspath{{\subfix{../imagenes/}}}
\begin{document}

\section{Sintesíntesis sonora}

\avisoesbozo

Sintesíntesis es un neologismo para referirse a la síntesis de sintetizadores. Este es un rubro llevado a cabo por muchos conglomerados e instituciones desde corporaciones (Roland, KORG), pasando por pequeñas empresas como Patch Point, Error Instruments, Mutable Instruments, hasta llegar a la producción artesanal unipersonal de sintetizadores en talleres de artistas.

En la siguiente sección se esbozan grandes influencias en el desarrollo de esta tesis, donde se selecciona a distintas luthieres y se explican los contextos en que trabajan.

\subsection{Peter Blasser}

Peter Blasser es un artista, luthier y empresario estadounidense, actualmente radicado en Berlín, Alemania.

Es el creador de familias de sintetizadores, conglomerados en familias o marcas, entre ellas destacando Ciat-Lonbarde, Shbobo, Ieaskul F. Mobenthey, y Tocante. Cada una tiene distintas filosofías y aproximaciones en cuanto a su fabricación, uso, afinación, interfaz.

Tras haber creado su taller de luthería en Estados Unidos, se asoció con Darrin Wiener, y trasladó y expandió sus operaciones a Alemania, donde en la tienda Patch Point ahora comercializa sus instrumentos y sigue desarrollando nuevas líneas. Tienen una tienda al público en Berlín, donde también se puede arrendar de forma gratuita un estudio para grabar y conocer los instrumentos. El taller de madera y cortes CNC está en Portugal, desacoplando la producción de los instrumentos en distintas capas de luthería.

\subsection{Émilie Gillet}

Émilie Gillet es una artista francesa, fundadora de la empresa Mutable Instruments \parencite{perfectcircuit_mutableinstruments_sinfecha}. Se destaca a Émilie en este contexto, ya que su obra está enfocada en el diseño integral de sus instrumentos, incluyendo interfaz de usuario, firmware, software, hardware y documentación.

Su enfoque podría haber sido cerrado y privado y haber sido igualmente un enorme aporte en el desarrollo de nuevas síntesis, pero además tiene un carácter popular, al haber publicado y abierto sus procesos, compartiendo sus herramientas de taller, y la documentación y los códigos que viven en sus sintetizadores. Esto ha permitido que otras fabricantes puedan también adaptar sus estrategias a nuevos módulos, incluyendo a producir de forma legal clones de sus instrumentos.

\newpage

\end{document}
